% Part: incompleteness
% Chapter: theories-computability
% Section: inseparability

\documentclass[../../../include/open-logic-section]{subfiles}

\begin{document}

\olfileid{inc}{tcp}{ins}

\olsection{$\Th{Q}$ وچ ثابت ہو سکّن والے تے قابلِ رد جملیاں دے
  سیٹ قابلِ حساب طریقے نال وکھ نہیں کیتے جا سکدے}

منّو $\Th{\bar Q}$ اوہناں جملیاں دا سیٹ اے جنہاں دیاں
{\em نفیاں} $\Th{Q}$ وچ ثابت ہو سکدیاں نیں، یعنی
$\Th{\bar Q} = \Setabs{!A}{\Th{Q} \Proves
  \lnot !A}$۔ ایہ قابلِ رد جملیاں دا سیٹ اے، نظریے دا پورا
متمم نہیں۔ یاد رکھو کہ بے سانجھے سیٹ $A$ تے $B$ قابلِ حساب
طریقے نال وکھ نہیں کیتے جا سکدے جے کوئی ایہو جیہا قابلِ حساب
سیٹ $C$ نہ ہووے جیہدے لئی $A
\subseteq C$ تے $B \subseteq \Complement{C}$۔

\begin{lem}
$\Th{Q}$ تے $\Th{\bar Q}$ قابلِ حساب طریقے نال وکھ نہیں کیتے
جا سکدے۔
\end{lem}

\begin{proof}
فرض کرو $C$ اک قابلِ حساب سیٹ اے جیہدے لئی $\Th{Q} \subseteq C$
تے $\Th{\bar Q} \subseteq \Complement{C}$۔ $R(x,y)$ نوں ایہ
تعلق منّ لئیے۔
\begin{quote}
$x$، فارمولے $!D(u)$ دا کوڈ اے تے $!D(\num y)$، $C$ وچ اے۔
\end{quote}
اسیں وکھاواں گے کہ $R(x,y)$ ہمہ گیر قابلِ حساب تعلق اے،
تے ایس توں تضاد نکلے گا۔ فارمولے دا کوڈ جانچنا، اوہدے وچ عددی
نام رکھنا تے دتے ہوئے قابلِ حساب سیٹ دی رکنیت جانچنا، تینوں کم
مؤثر طریقے نال ہو سکدے نیں۔ ایس لئی پورا تعلق قابلِ حساب اے۔
جے پہلا انپٹ فارمولے دا صحیح کوڈ نہ ہووے، تاں تعلق قائم نہیں ہوندا۔

فرض کرو $S(y)$ قابلِ حساب اے تے $!D_S(u)$، $\Th{Q}$ وچ
اوہدی نمائندگی کردا اے۔ فیر
\begin{eqnarray*}
S(\num n) & \lif & \Th{Q} \Proves !D_S(\num n) \\
& \lif & !D_S(\num n) \in C
\end{eqnarray*}
تے
\begin{eqnarray*}
\lnot S(\num n) & \lif & \Th{Q} \Proves \lnot !D_S(\num n) \\
& \lif & !D_S(\num n) \in \Th{\bar Q} \\
& \lif & !D_S(\num n) \not\in C
\end{eqnarray*}
ایس لئی $S(y)$ تے $R(\#(!D_S(u)),y)$ ہم ارز نیں۔ کوڈ آزاد
متغیر والے اصل فارمولے دا اے۔ ہر انپٹ لئی اوہدی عددی مثال
سیٹ دی رکنیت جانچن لئی ورتی جاندی اے۔ ایس طرح ہر اک رکنی
قابلِ حساب تعلق لئی اک مقرر قطار ملدی اے۔ فیر ہمہ گیر تعلق
فرض کرن نال پچھلا قطری تضاد نکلدا اے۔
% PNB-CMP-068: اصل آخری کوڈ وچ عددی مثال دی بجائے آزاد متغیر والا فارمولا؛واحد محدود ماخذی تصحیح۔
\end{proof}

\end{document}
